\section{関連研究}\label{sec:related}

\subsection{差分プライバシー}

Dwork ら~\cite{dwork2006calibrating}は差分プライバシーの概念を導入し、ラプラスメカニズムによるカウントクエリのプライバシー保護を提案した。
$(\varepsilon, \delta)$-差分プライバシーの緩和定義~\cite{dwork2014algorithmic}はガウスメカニズムを可能にし、合成定理（逐次合成・高度合成）~\cite{dwork2010boosting}は複数クエリにわたるプライバシ保証の管理を可能にした。
Sparse Vector Technique (SVT)~\cite{dwork2009complexity,hardt2010simple}は、閾値判定のみが必要な場合にプライバシ予算を節約する技法である。

\subsection{差分プライバシーと頻出パターンマイニング}

頻出アイテムセットマイニングへの差分プライバシーの適用は複数の研究で取り組まれている。
Zeng ら~\cite{zeng2012differentially}はトランザクションの切断（truncation）とFP-tree への指数メカニズムの適用を提案した。
Li ら~\cite{li2012privbasis}はPrivBasisアルゴリズムを提案し、基底集合からのノイズ付きサポート推定を行った。
Bhaskar ら~\cite{bhaskar2010discovering}は指数メカニズムを用いた top-$k$ 頻出アイテムセットの発見を提案した。
これらの研究はいずれもグローバルなサポート値を対象としており、時間的な密集パターンは扱っていない。

\subsection{密集区間マイニング}

スライディングウィンドウに基づく密集区間検出~\cite{suzuki2024dense}は、Apriori性質を利用してアイテムセットの時間的集中を効率的に検出する。
ウィンドウサイズ$W$と最小サポート$\theta$をパラメータとし、ウィンドウ内の共起カウントが$\theta$以上となる連続区間を出力する。
本手法は時系列マイニングとトランザクション分析の交差領域に位置するが、プライバシー保証は考慮されていない。

\subsection{差分プライバシーと時系列・ストリームマイニング}

時系列データへの差分プライバシーの適用も研究されている。
Chan ら~\cite{chan2011private}はストリームデータに対する連続的なカウントクエリの差分プライバシーを研究した。
Dwork ら~\cite{dwork2010differential}は連続観測モデルにおけるプライバシ保証を分析した。
Kellaris ら~\cite{kellaris2014differentially}は時系列データに対するプライバシー保護を議論した。
しかし、これらは連続的なカウント発行に焦点を当てており、密集区間の検出というパターンマイニング的側面は扱っていない。

\subsection{本研究の位置づけ}

以上の関連研究を踏まえ、本研究は「密集区間マイニング」と「差分プライバシー」の交差領域における空白を埋めるものである。
具体的には、ウィンドウカウントクエリの感度解析、閾値安定性の理論的保証、およびApriori型多レベル探索における予算合成を統合したフレームワークを提案する。
